\textbf{Commun à tous les candidats}

\medskip

Dans cet exercice, les résultats des probabilités demandées seront, si nécessaire, arrondis au millième.


























\medskip

\begin{enumerate}
\item 
	\begin{enumerate}
		\item On traduit la situation par un arbre pondéré.
		
\begin{center}
\pstree[treemode=R,nodesepB=3pt,levelsep=2.75cm]{\TR{}}
{\pstree{\TR{$M$~~}\taput{0,4}}
	{\TR{$T$}\taput{0,9}
	\TR{$\overline{T}$}\tbput{0,1}
	}
\pstree{\TR{$\overline{M}$~~}\tbput{0,6}}
	{\TR{$T$}\taput{0,15}
	\TR{$\overline{T}$}\tbput{0,85}
	}	 
}		
\end{center}

\medskip
	
		\item 
Il faut trouver $P(M \cap T) = P(M) \times P_M(T) = 0,4 \times 0,9 = 0,36$.
		\item 
On a de même $P\left(\overline{M} \cap T\right) = P\left(\overline{M}\right) \times P_{\overline{M}}(T) = 0,6 \times  0,15 = 0,09$.

D'après la loi des probabilités totales :

$P(T) = P(M \cap T) + P\left(\overline{M} \cap T\right) = 0,36 + 0,09 = 0,45$.
		\item 
		Il faut trouver $P_T(M) = \dfrac{P(M \cap T)}{P(T)} = \dfrac{0,36}{0,45} = \dfrac{36}{45} = \dfrac{9 \times 4}{9 \times 5} = \dfrac{4}{5} = \dfrac{8}{10} = 0,8$.
	\end{enumerate}
\item 


	\begin{enumerate}
		\item 
On suppose que le nombre de chats est assez important pour que l'on puisse assimiler le choix des 20 chats à un tirage avec remise.
		
La variable $X$ suit donc une loi binomiale de paramètres $n = 20$ et de probabilité $p = 0,45$ trouvé à la question \textbf{1. c.}.
		\item 
On a $p\left(X = 5\right) = \binom{20}{5} \times 0,45^5 \times (1 - 0,45)^{20-5} = \np{15504}\times 0,45^5  \times 0,55^{15} \approx \np{0,0365}$ soit environ 0,037.

		\item 
La calculatrice donne $P(X < 9) \approx 0,414$.
		\item 
On sait que l'espérance $E = n \times p = 20 \times 0,45 = 9$.
		
Cela signifie que sur un grand nombre d'échantillons il y aura en moyenne 9 chats positifs par échantillon de 20.
	\end{enumerate}		
\item 
	\begin{enumerate}
		\item 
On a encore une loi binomiale de paramètres $n$ et de probabilité d'être positif de $0,45$.

On a $P(X = 0) = \binom{n}{0}\times 0,45^0 \times 0,55^n = 0,55^n$.

Donc $p_n = 1 - P(X = 0) =  1 - 0,55^n$.
		\item












En partant de $n=0$, le programme calcule $p_n$ et augmente la taille de l'échantillon de 1 tant que $p_n < 0,99$.

		\item 
		On cherche donc $n$ tel que $1 - 0,55^n \geqslant 0,99 \iff 0,01 \geqslant 0,55^n$, d'où par croissance de la fonction logarithme népérien : $\ln 0,01 \geqslant n \ln 0,55 \iff \dfrac{\ln 0,01}{\ln 0,55} \leqslant n$ (car $\ln 001 < 0$).
		
		Or $\dfrac{\ln 0,01}{\ln 0,55} \approx 7,7$.
		
Conclusion : le programme renvoie la valeur 8.
	\end{enumerate}
\end{enumerate}

\bigskip
